\begin{proposition}[Minimal quotient for the SHIFT/Lugano witness]
Let \(L\) denote the eight-valued outcome of the SHIFT measurement and let
\(C=f(L)\) denote the four-valued Lugano output class obtained by applying
\(f(0)=f(1)=0\) and \(f(+)=f(-)=1\) componentwise. For uniformly distributed
computational inputs, \(C\) is uniformly distributed over four values and equals
the Lugano target \(w_L(a,b,c)\) with probability one.

Let \(Y\) be any coarse-grained outcome obtained from \(L\) by an arbitrary
stochastic channel \(P(Y|L)\), with \(|Y|\le m\). Then the optimal decoding
success probability satisfies
\[
    P_{\rm succ}\le \frac{\min(m,4)}{4}.
\]
In particular, any coarse-graining with at most three outcomes cannot violate
the Lugano causal bound \(3/4\), whereas the four-outcome quotient
\(Y=C=f(L)\) achieves \(P_{\rm succ}=1\).
\end{proposition}

\begin{proof}
For a decoder \(d:Y\to\{1,2,3,4\}\),
\[
P_{\rm succ}=\sum_y P(Y=y,C=d(y)).
\]
The optimal decoder chooses, for each \(y\), a class \(c_y\) maximizing
\(P(Y=y,C=c)\), so
\[
P_{\rm succ}=\sum_y \max_c P(Y=y,C=c)
            =\sum_y P(Y=y,C=c_y).
\]
Grouping terms by \(c_y\),
\[
P_{\rm succ}
 = \sum_c \sum_{y:c_y=c} P(Y=y,C=c)
 \le \sum_{c\in\{c_y:y\in Y\}} P(C=c).
\]
At most \(m\) classes appear in this last sum, and \(P(C=c)=1/4\) for each of
the four Lugano classes. Therefore \(P_{\rm succ}\le m/4\) for \(m\le 4\),
and \(P_{\rm succ}\le 1\) otherwise. The bound is achieved by the quotient
measurement \(Y=f(L)\), which has four outcomes and preserves the Lugano
target exactly.
\end{proof}